\section{Reduction to controlled predecessors and complete splitting}
\label{sec:reduction}

For a prime $p\nmid a$, failure of $a$ to generate $\mathbb F_p^\times$
means that some prime divides the index
\[
 i_p(a)=\frac{p-1}{\ord_p(a)}.
\]
We construct primes for which the possible divisors of this index belong
to two ranges. A large divisor forces $p$ to divide one of a short list
of integers $a^j-1$; the remaining divisors are controlled by complete
splitting in Kummer fields. This is the index-divisor reduction underlying
Hooley's work \cite[Sections~4--7]{Hooley67}.

The first proposition supplies the primes, with a fixed congruence that
will force quadratic nonresiduosity.

\begin{proposition}[Primes with controlled predecessors]
\label{prop:controlled}\label{prop:prime-reservoir}
Let $M$ be a fixed positive multiple of $8$, let $c\in\{2,4\}$, and let
$u$ be an integer satisfying
\begin{equation}
 (u,M)=1,\qquad c\mid u-1,\qquad
 \left(\frac{u-1}{c},\frac Mc\right)=1.
 \label{alg:progression-conditions}
\end{equation}
For every sufficiently large real $x$, there are
$\gg_{M,c,u}x/L^2$ distinct primes $p\in(x,2x)$ satisfying
\begin{equation}
 p\equiv u\pmod M,\qquad p-1=crQ,
 \qquad Q>x^{0.9}\ \text{prime},
 \label{eq:controlled-predecessor}
\end{equation}
where $r$ is a positive integer all of whose prime divisors lie in
\[
 \bigl(\exp(L^{0.1}),\exp(L^{0.3})\bigr).
\]
\end{proposition}

We prove Proposition~\ref{prop:controlled} in
Section~\ref{sec:construction}. For the second input, write $\mu_m$ for
the group of $m$th roots of unity. For a nonzero integer $a$ and a prime
$q$, put
\begin{equation}
 E_q=\Q(\mu_q),\qquad K_q=\Q(\mu_q,a^{1/q}).
 \label{alg:kummer-fields}
\end{equation}
The choice of the root $a^{1/q}$ does not affect $K_q$.

\begin{proposition}[Uniform bound for complete splitting]
\label{prop:splitting}
Fix an integer $a$ with $|a|>1$, and put $\delta_0=10^{-6}$.
For all sufficiently large real $x$ and all primes $q$ sufficiently
large in terms of $a$, with $q\le\exp(L^{0.3})$, one has
\begin{equation}
 \#\{p\text{ prime}:x<p\le2x,\ p\text{ splits completely in }K_q\}
 \ll_a \frac{x}{q(q-1)L}+x^{1-\delta_0}.
 \label{eq:uniform-splitting}
\end{equation}
The implied constant and the threshold for $x$ are independent of $q$.
For $a=2$, every prime $q\ge5$ is allowed, with absolute constants and
an absolute threshold for $x$.
\end{proposition}

Section~\ref{sec:splitting} proves this proposition from the Hecke
zero-free theorem, Theorem~\ref{thm:hecke-zero-free}. To apply it, we first record
the elementary interpretation of complete splitting.

\begin{lemma}[Splitting test]\label{alg:splitting-test}
Let $a\ne0$ be an integer, and let $p,q$ be primes with $p\nmid aq$.
Then $p$ splits completely in $K_q$ if and only if
\begin{equation}
 p\equiv1\pmod q,\qquad a^{(p-1)/q}\equiv1\pmod p.
 \label{eq:kummer-frobenius-test}
\end{equation}
\end{lemma}

\begin{proof}
The field $K_q$ is the splitting field of $T^q-a$. Fix a prime
$\mathfrak P$ of $K_q$ above $p$. The roots are integral and have distinct
reductions modulo $\mathfrak P$, since $p$ does not divide the polynomial
discriminant $\pm q^q a^{q-1}$. Every inertia element permutes the roots
and fixes their reductions, so it fixes every root. Thus $p$ is
unramified. Let $\sigma$ be its arithmetic Frobenius at $\mathfrak P$.
On $\mu_q$, this automorphism acts by $\zeta\mapsto\zeta^p$, so it fixes
$E_q$ exactly when $p\equiv1\pmod q$. Under this condition, with
$\alpha=a^{1/q}$, we have
\[
 \frac{\sigma(\alpha)}{\alpha}\in\mu_q,
 \qquad
 \frac{\sigma(\alpha)}{\alpha}
 \equiv\alpha^{p-1}=a^{(p-1)/q}\pmod{\mathfrak P}.
\]
The division is valid because $p\nmid a$. Reduction on $\mu_q$ is
injective at $p\ne q$, so the second congruence in
\eqref{eq:kummer-frobenius-test} is equivalent to $\sigma(\alpha)=\alpha$.
Together the two conditions say that $\sigma$ fixes the generators of
$K_q$. Since $K_q/\Q$ is Galois, this is precisely complete splitting.
\end{proof}

We now choose a progression of quadratic nonresidues satisfying all
the hypotheses of Proposition~\ref{prop:controlled}.

\begin{lemma}[A progression of quadratic nonresidues]
\label{alg:quadratic-progression}
Let $a\in\Z$ be neither $-1$ nor an integer square, and set
\[
 M=8\prod_{\substack{\ell\mid a\\\ell\text{ odd prime}}}\ell.
\]
There are $c\in\{2,4\}$ and an integer $u$ satisfying
\eqref{alg:progression-conditions} such that every prime $p\equiv u\pmod M$
satisfies $p\nmid a$ and $(a/p)=-1$.
\end{lemma}

\begin{proof}
Write $a=d_0t^2$, where $t\ge1$ is an integer and $d_0\ne1$ is signed
squarefree. We choose $u$ to be a unit modulo $M$, with
$u\not\equiv1\pmod\ell$ for every odd prime $\ell\mid M$, and
$u\bmod8\in\{3,5,7\}$. These restrictions already imply $p\nmid a$.

Suppose first that an odd prime $j\ge5$ divides $d_0$. Fix any of the
allowed residues modulo $8$ and at the odd primes other than $j$.
Quadratic reciprocity expresses $(d_0/p)$ as a fixed sign times $(p/j)$.
Among the nonzero residues modulo $j$ other than $1$, there are
$(j-3)/2\ge1$ squares and $(j-1)/2$ nonsquares. We can therefore choose
$u\bmod j$ so that $(d_0/p)=-1$.

If no such $j$ exists, the seven possibilities for $d_0$ are treated by
the following table. Whenever $3\mid M$, also take $u\equiv2\pmod3$.
\begin{center}
\begin{tabular}{c|rrrrrrr}
$d_0$ & $-1$ & $2$ & $-2$ & $3$ & $-3$ & $6$ & $-6$\\ \hline
$u\bmod8$ & $3$ & $3$ & $7$ & $5$ & $5$ & $3$ & $7$
\end{tabular}
\end{center}
The supplementary laws give the first three entries. For the last four,
the combined residues modulo $24$ are respectively $5,5,11,23$;
quadratic reciprocity gives $(d_0/p)=-1$ in each case. Odd primes
dividing $M$ but not $d_0$ occur only in the square factor $t^2$ and can
be assigned any allowed residue. The Chinese remainder theorem combines
these choices into a class $u\pmod M$.

If $u\equiv3$ or $7\pmod8$, take $c=2$; if $u\equiv5\pmod8$, take
$c=4$. In each case $(u-1)/c$ is odd. For every odd $\ell\mid M$,
the condition $u\not\equiv1\pmod\ell$ also makes $(u-1)/c$ nonzero
modulo $\ell$. This proves \eqref{alg:progression-conditions}.
Finally, $(a/p)=(d_0/p)=-1$, since $p\nmid t$.
\end{proof}

\begin{proof}[Proof of Theorem~\ref{thm:primitive-root}]
Fix $a$ as in the theorem, so $|a|>1$. Choose $M,c,u$ by
Lemma~\ref{alg:quadratic-progression}, and take the
$\gg_a x/L^2$ primes supplied by Proposition~\ref{prop:controlled}.
For each such prime $p$, Euler's criterion gives
$a^{(p-1)/2}\equiv-1\pmod p$. Consequently $i_p(a)$ is odd: an even
index would make $\ord_p(a)$ divide $(p-1)/2$.

If $i_p(a)>1$, choose a prime divisor $q$ of this index. Since
$p-1=crQ$ and $c$ is a power of $2$, we have $q\mid rQ$. Moreover
$\ord_p(a)\mid(p-1)/q$, so \eqref{eq:kummer-frobenius-test} holds.
We bound the number of such primes $p$ in the two possible ranges of $q$.

If $q=Q$, then
\[
 1\le\frac{p-1}{Q}<2x^{0.1}.
\]
Thus $p$ divides the positive integer
\[
 B_a(x)=\prod_{1\le j\le2x^{0.1}}|a^j-1|.
\]
No factor vanishes because $|a|>1$, and
\[
 \log B_a(x)
 \le\sum_{1\le j\le2x^{0.1}}\bigl(j\log|a|+\log2\bigr)
 \ll_a x^{0.2}.
\]
Every distinct prime divisor $p>x$ contributes at least $L$ to this
logarithm. The number of primes lost in this case is therefore
$O_a(x^{0.2}/L)=o(x/L^2)$.

If $q\mid r$, then
$\exp(L^{0.1})<q<\exp(L^{0.3})$. For large $x$, all these primes are
sufficiently large for Proposition~\ref{prop:splitting}; also $p\nmid aq$.
Lemma~\ref{alg:splitting-test} and a union bound give at most
\begin{align}
 &\ll_a \frac{x}{L}
       \sum_{m>\exp(L^{0.1})}\frac1{m(m-1)}
       +x^{1-\delta_0}\exp(L^{0.3})\notag\\
 &\ll_a \frac{x}{L}\exp(-L^{0.1})
       +x\exp(-\delta_0L+L^{0.3})
       =o(x/L^2)
 \label{alg:exceptional-small-divisors}
\end{align}
primes, where the first sum is over integers. The first estimate uses
\eqref{eq:uniform-splitting} for each possible $q$; the second uses the
telescoping sum of $1/(m(m-1))$.

Subtracting these two losses leaves $\gg_a x/L^2$ primes with
$i_p(a)=1$. All choices depend only on $a$, and both estimates hold for
every sufficiently large real $x$. This proves the theorem.
\end{proof}

For $a=2$, the theorem gives $\gg x/L^2$ primes in $(x,2x)$ with
primitive root $2$, with absolute constants. This case can use the
particularly simple choice $(M,c,u)=(8,4,5)$ in
Proposition~\ref{prop:controlled}: the predecessors are $4rQ$, and
$p\equiv5\pmod8$ gives $(2/p)=-1$ directly.

\begin{remark}
Negative squares are covered by the row $d_0=-1$ in the progression
lemma. Odd perfect powers also require no separate exclusion. If
$a=b^h$ with fixed odd $h>1$, every prime divisor of $h$ is eventually
smaller than every prime factor of $rQ$; hence $(h,p-1)=1$ for the
constructed primes. This explains directly how the controlled
predecessor removes the obstruction from the fixed exponent.
\end{remark}
